% GENERATED FILE. Edit canonical lessons, then run npm run generate:books.
% canonical-source-hash: fnv1a64:93bf943da793ed04
% canonical-lessons: GE-C411-fahrkartenautomat, GE-C411-fahrschein, GE-C411-fahrpreis, GE-C411-entwerten, GE-C411-loesen

\chapter{Fahrkarten kaufen}
\label{ch:ge-fahrkarten-kaufen}
\hlchaptermodality{411}

\begin{chapteropening}
\textbf{By the end of this chapter:} \emph{I can buy a ticket: the ticket machine, the travel ticket, the fare, stamping a ticket and getting one.}

Complete the last lesson of chapter 411: I can buy a ticket: the ticket machine, the travel ticket, the fare, stamping a ticket and getting one.
\end{chapteropening}

\section[der Fahrkartenautomat]{der Fahrkartenautomat — the ticket machine}

\label{lesson:GE-C411-fahrkartenautomat}



\begin{quote}
Before the new one: say the German for the pedestrian, then the German for the cyclist.
\end{quote}

\subsection*{You'll want to know: der Fahrkartenautomat}
\textbf{der Fahrkartenautomat} — \textquotedblleft{}the ticket machine\textquotedblright{}.

\begin{quote}
Der Fahrkartenautomat nimmt auch Kreditkarten.
\end{quote}

\begin{quote}
\emph{The ticket machine also takes credit cards.}
\end{quote}

Plural: \emph{die Fahrkartenautomaten}; weak noun: \emph{den Fahrkartenautomaten}.

\begin{cousinweb}
\emph{Fahrkartenautomat} is built from \emph{Fahrkarte} 'ticket', from \emph{fahren} and \emph{Karte} 'card' (from Latin \emph{charta}, Lewis \& Short), and \emph{Automat} 'machine', which goes back to Greek \emph{autómatos} 'acting by itself' (EtymDB).
\end{cousinweb}

\begin{grammarlens}[title={What you've built}]
One more word for telling, explaining and describing at B1.
\end{grammarlens}

\subsection*{Your turn}
\begin{itemize}
\raggedright
  \item \emph{Say it:} \emph{der Fahrkartenautomat}
  \item \emph{Say it:} \emph{der Fahrkartenautomat}, once more
  \item \emph{Say it:} a sentence of your own with \emph{der Fahrkartenautomat}
  \item \emph{Recall:} say the German for the pedestrian, then the German for the cyclist, then say \emph{der Fahrkartenautomat} again
\end{itemize}

\begin{tcolorbox}[breakable,colback=teal!4,colframe=teal!35!black,title={Before you move on}]
What does der Fahrkartenautomat mean? (\textquotedblleft{}The ticket machine\textquotedblright{}.) Say it once more.
\end{tcolorbox}
\section[der Fahrschein]{der Fahrschein — the travel ticket (formal)}

\label{lesson:GE-C411-fahrschein}



\begin{quote}
Before the new one: say the German for the cyclist, then the German for the ticket machine.
\end{quote}

\subsection*{You'll want to know: der Fahrschein}
\textbf{der Fahrschein} — \textquotedblleft{}the travel ticket (formal)\textquotedblright{}.

\begin{quote}
Ohne Fahrschein darf man nicht in die Straßenbahn.
\end{quote}

\begin{quote}
\emph{You may not get on the tram without a ticket.}
\end{quote}

Plural: \emph{die Fahrscheine}. The word on official signs; people mostly say \emph{die Fahrkarte} or \emph{das Ticket}.

\begin{cousinweb}
\emph{Fahrschein} is built from \emph{fahren} 'to travel' and \emph{Schein} 'certificate, slip', both inherited: \emph{Schein} goes back to Proto-Germanic \emph{skīnaną} 'to shine' (EtymDB).
\end{cousinweb}

\begin{grammarlens}[title={What you've built}]
One more word for telling, explaining and describing at B1.
\end{grammarlens}

\subsection*{Your turn}
\begin{itemize}
\raggedright
  \item \emph{Say it:} \emph{der Fahrschein}
  \item \emph{Say it:} \emph{der Fahrschein}, once more
  \item \emph{Say it:} a sentence of your own with \emph{der Fahrschein}
  \item \emph{Recall:} say the German for the cyclist, then the German for the ticket machine, then say \emph{der Fahrschein} again
\end{itemize}

\begin{tcolorbox}[breakable,colback=teal!4,colframe=teal!35!black,title={Before you move on}]
What does der Fahrschein mean? (\textquotedblleft{}The travel ticket (formal)\textquotedblright{}.) Say it once more.
\end{tcolorbox}
\section[der Fahrpreis]{der Fahrpreis — the fare}

\label{lesson:GE-C411-fahrpreis}



\begin{quote}
Before the new one: say the German for the ticket machine, then the German for the travel ticket.
\end{quote}

\subsection*{You'll want to know: der Fahrpreis}
\textbf{der Fahrpreis} — \textquotedblleft{}the fare\textquotedblright{}.

\begin{quote}
Für Kinder ist der Fahrpreis halb so hoch.
\end{quote}

\begin{quote}
\emph{For children the fare is half as much.}
\end{quote}

Plural: \emph{die Fahrpreise}.

\begin{cousinweb}
\emph{Fahrpreis} is built from \emph{fahren} 'to travel' and \emph{Preis} 'price', which is usually traced through Old French \emph{pris} to Latin \emph{pretium} 'price, value' (Lewis \& Short).
\end{cousinweb}

\begin{grammarlens}[title={What you've built}]
One more word for telling, explaining and describing at B1.
\end{grammarlens}

\subsection*{Your turn}
\begin{itemize}
\raggedright
  \item \emph{Say it:} \emph{der Fahrpreis}
  \item \emph{Say it:} \emph{der Fahrpreis}, once more
  \item \emph{Say it:} a sentence of your own with \emph{der Fahrpreis}
  \item \emph{Recall:} say the German for the ticket machine, then the German for the travel ticket, then say \emph{der Fahrpreis} again
\end{itemize}

\begin{tcolorbox}[breakable,colback=teal!4,colframe=teal!35!black,title={Before you move on}]
What does der Fahrpreis mean? (\textquotedblleft{}The fare\textquotedblright{}.) Say it once more.
\end{tcolorbox}
\section[entwerten]{entwerten — to stamp, to validate (a ticket)}

\label{lesson:GE-C411-entwerten}



\begin{quote}
Before the new one: say the German for the travel ticket, then the German for the fare.
\end{quote}

\subsection*{You'll want to know: entwerten}
\textbf{entwerten} — \textquotedblleft{}to stamp, to validate (a ticket)\textquotedblright{}.

\begin{quote}
Im Bus entwertet sie ihre Fahrkarte.
\end{quote}

\begin{quote}
\emph{On the bus she stamps her ticket.}
\end{quote}

Literally to take the value away: once stamped, the ticket cannot be used again. Inseparable: \emph{sie hat entwertet}.

\begin{cousinweb}
\emph{entwerten} is built from \emph{ent-} 'away, un-' and \emph{Wert} 'value, worth', an inherited Germanic word that is a cousin of English \emph{worth}.
\end{cousinweb}

\begin{grammarlens}[title={What you've built}]
One more word for telling, explaining and describing at B1.
\end{grammarlens}

\subsection*{Your turn}
\begin{itemize}
\raggedright
  \item \emph{Say it:} \emph{entwerten}
  \item \emph{Say it:} \emph{entwerten}, once more
  \item \emph{Say it:} a sentence of your own with \emph{entwerten}
  \item \emph{Recall:} say the German for the travel ticket, then the German for the fare, then say \emph{entwerten} again
\end{itemize}

\begin{tcolorbox}[breakable,colback=teal!4,colframe=teal!35!black,title={Before you move on}]
What does entwerten mean? (\textquotedblleft{}To stamp, to validate (a ticket)\textquotedblright{}.) Say it once more.
\end{tcolorbox}
\section[lösen]{lösen — to get a ticket; to solve; to loosen}

\label{lesson:GE-C411-loesen}



\begin{quote}
Before the new one: say the German for the fare, then the German for to stamp.
\end{quote}

\subsection*{You'll want to know: lösen}
\textbf{lösen} — \textquotedblleft{}to get a ticket; to solve; to loosen\textquotedblright{}.

\begin{quote}
Am Bahnhof löst er eine Fahrkarte nach Hause.
\end{quote}

\begin{quote}
\emph{At the station he gets a ticket home.}
\end{quote}

\emph{Eine Fahrkarte lösen} sounds a little formal. \emph{Ein Problem lösen} — to solve a problem.

\begin{cousinweb}
\emph{lösen} is inherited through Old High German \emph{lōsen} from Proto-Germanic \emph{lausijaną}, built on \emph{lausaz} 'loose, free' (EtymDB). English \emph{loose} is a cousin.
\end{cousinweb}

\begin{grammarlens}[title={What you've built}]
One more word for telling, explaining and describing at B1.
\end{grammarlens}

\subsection*{Your turn}
\begin{itemize}
\raggedright
  \item \emph{Say it:} \emph{lösen}
  \item \emph{Say it:} \emph{lösen}, once more
  \item \emph{Say it:} a sentence of your own with \emph{lösen}
  \item \emph{Recall:} say the German for the fare, then the German for to stamp, then say \emph{lösen} again
\end{itemize}

\begin{tcolorbox}[breakable,colback=teal!4,colframe=teal!35!black,title={Before you move on}]
What does lösen mean? (\textquotedblleft{}To get a ticket; to solve; to loosen\textquotedblright{}.) Say it once more.
\end{tcolorbox}
